\section{Fixed-determinant projective Plancherel}
\label{sec:tpc33-projective-plancherel}

We pass to a finite model in which the two-parameter affine coefficient
family is complete.  The terminology ``projective Plancherel'' is
descriptive: the proof is finite Fourier orthogonality in the coordinates
\(b=a/s\) and \(c=h/s\).  Thus the second moment averages both the
projective ratio \(b\) and the nonzero affine offset \(c\); it gives no
almost-every-\(b\) statement by itself.  The useful feature here is its
specialization to one nonzero determinant fiber and the explicit defect
left by homogeneous dilations.

Let \(q\) be prime, let \(h\in\F_q^\times\), and let
\(f,g:\F_q\to\C\).  We use the unnormalized conventions
\begin{align}
 \bar f&=\frac1q\sum_{x\in\F_q}f(x),
 &f^\circ&=f-\bar f,
 \label{eq:tpc33-finite-centering}\\
 \|f\|_2^2&=\sum_{x\in\F_q}|f(x)|^2,
 &\widehat f(r)&=\sum_{x\in\F_q}f(x)e_q(-rx),
 \label{eq:tpc33-finite-normalization}
\end{align}
where \(e_q(z)=\exp(2\pi iz/q)\).  For
\(a,s\in\F_q^\times\), define
\begin{equation}
 \cC_{a,s}^{(h)}(f,g)
 :=\sum_{D\in\F_q}
 f(D)g\!\left(s^{-1}(aD+h)\right).
 \label{eq:tpc33-projective-correlation}
\end{equation}
This is the correlation along the affine line \(sU-aD=h\).

\begin{theorem}[Exact determinant-fiber second moment]
\label{thm:tpc33-projective-plancherel}
With the preceding conventions,
\begin{align}
 &\boxed{
 \sum_{a,s\in\F_q^\times}
 \left|
  \cC_{a,s}^{(h)}(f,g)-q\bar f\bar g
 \right|^2}
 \notag\\[-0.2em]
 &\boxed{\qquad
 =q\|f^\circ\|_2^2\|g^\circ\|_2^2
 -\sum_{b\in\F_q^\times}
  \left|\sum_{D\in\F_q}f^\circ(D)g^\circ(bD)\right|^2.}
 \label{eq:tpc33-projective-plancherel}
\end{align}
In particular,
\begin{equation}
 \boxed{
 \sum_{a,s\in\F_q^\times}
 \left|
  \cC_{a,s}^{(h)}(f,g)-q\bar f\bar g
 \right|^2
 \le q\|f^\circ\|_2^2\|g^\circ\|_2^2.}
 \label{eq:tpc33-projective-upper-bound}
\end{equation}
\end{theorem}

\begin{proof}
Put
\begin{equation}
 b=a/s,\qquad c=h/s.
 \label{eq:tpc33-projective-coordinates}
\end{equation}
Because \(h\ne0\), the map
\((a,s)\mapsto(b,c)\) is a bijection of
\((\F_q^\times)^2\): its inverse is
\(s=h/c\), \(a=bs\).  Centering both functions gives
\begin{equation}
 \cC_{a,s}^{(h)}(f,g)-q\bar f\bar g
 =R_{b,c},
 \qquad
 R_{b,c}:=\sum_Df^\circ(D)g^\circ(bD+c).
 \label{eq:tpc33-centered-affine-correlation}
\end{equation}
The mixed term containing \(g^\circ\) vanishes because an affine permutation
preserves its zero sum, while the other vanishes because
\(\sum_D f^\circ(D)=0\).

Temporarily include \(c=0\).  Fourier transformation in \(c\) gives
\begin{equation}
 \widehat R_b(r)
 :=\sum_{c\in\F_q}R_{b,c}e_q(-rc)
 =\widehat f^\circ(-br)\widehat g^\circ(r).
 \label{eq:tpc33-R-fourier}
\end{equation}
Parseval and \(\widehat f^\circ(0)=\widehat g^\circ(0)=0\) imply
\begin{align}
 \sum_{b\ne0}\sum_{c\in\F_q}|R_{b,c}|^2
 &=\frac1q\sum_{r\ne0}|\widehat g^\circ(r)|^2
   \sum_{b\ne0}|\widehat f^\circ(-br)|^2
 \notag\\
 &=\frac1q
   \left(\sum_{\rho\ne0}|\widehat f^\circ(\rho)|^2\right)
   \left(\sum_{r\ne0}|\widehat g^\circ(r)|^2\right)
 \notag\\
 &=q\|f^\circ\|_2^2\|g^\circ\|_2^2.
 \label{eq:tpc33-complete-affine-energy}
\end{align}
The determinant-\(h\) family has \(c\ne0\).  Subtracting its omitted
\(c=0\) energy,
\[
 \sum_{b\ne0}|R_{b,0}|^2
 =\sum_{b\ne0}
   \left|\sum_Df^\circ(D)g^\circ(bD)\right|^2,
\]
proves \eqref{eq:tpc33-projective-plancherel}.  Dropping this nonnegative
term proves \eqref{eq:tpc33-projective-upper-bound}.
\end{proof}

\begin{remark}[The dilation-alignment defect]
The subtracted term is the complete energy of the homogeneous lines
\(U=bD\).  Thus the identity can be strictly smaller than its upper bound
when \(f^\circ\) and a multiplicative dilation of \(g^\circ\) align.  No
randomness assumption is made.
\end{remark}

\begin{corollary}[Exceptional determinant-fiber pairs]
\label{cor:tpc33-exceptional-projective-directions}
For every \(\Lambda>0\),
\begin{align}
 &\#\left\{(a,s)\in(\F_q^\times)^2:
 \left|\cC_{a,s}^{(h)}(f,g)-q\bar f\bar g\right|>\Lambda
 \right\}
 \notag\\
 &\qquad\le
 \frac{q\|f^\circ\|_2^2\|g^\circ\|_2^2}{\Lambda^2}.
 \label{eq:tpc33-projective-exception-count}
\end{align}
\end{corollary}

\begin{proof}
Apply Chebyshev's inequality to
\eqref{eq:tpc33-projective-upper-bound}.
\end{proof}

The corollary is an average over the two-dimensional coefficient family
\((a,s)\), equivalently over \((b,c)\).  It gives no pointwise cancellation
for one prescribed physical pair \((a,s)\), nor an almost-every-\(b\)
conclusion, unless an additional argument controls the remaining coordinate.
